\section{Rational functions on the required grids}\label{sec:construction}

\subsection{Zeros, poles and reflection}

The $b$th form will be an alternating sum over the points $j/b$ with
$j$ odd. We need numerator zeros at the initial points of every such
grid, both to remove their contribution to the sum and to control
the denominators of its rational constant. The signs on odd summation
indices repeat after a shift by $4$, so poles at $0,-4,\ldots,-4n$
allow each partial-fraction sum to be shifted back to the beta series.
The distinct fraction classes that occur for $1\le b\le B$ are
\[
 \Theta=\left\{\frac ad:1\le d\le B,\ 1\le a\le d,\ (a,d)=1\right\}.
\]
The integer class is represented by $1$. Thus $|\Theta|=D_B$.
For example, at $B=3$ the classes are $1,1/2,1/3,2/3$.

Fix an integer $r\ge3$. If the rightmost zero at index $n$ is $rn$,
reflection about $-2n$, the midpoint of the pole interval, puts the
leftmost zero at $-(r+4)n$. The zero interval therefore has length
$(2r+4)n$. Abbreviate
\[
 D=D_B,\qquad L=2r+4,\qquad S=s+1,
 \qquad S>DL,
\]
where $s$ is even. The integer $r$ controls the length of the initial
zero interval. To normalize the denominators of each fraction class,
put
\begin{equation}\label{eq:K}
 A_d=d^L\prod_{p\mid d}p^{\lceil L/(p-1)\rceil},\qquad
 K=\prod_{\theta\in\Theta}A_{\operatorname{den}(\theta)},
\end{equation}
where the first product is over primes and
$\operatorname{den}(\theta)$ is the reduced denominator. The
progression estimate in Section~\ref{sec:arithmetic} explains this
normalization.

For every positive odd integer $n$, define
\begin{align}
 P_n(t)&=(t-rn)\prod_{\theta\in\Theta}
                    \prod_{j=0}^{Ln-1}(t-rn+j+\theta),\label{eq:P}\\
 C_n&=4^{Sn}K^n\frac{((n+1)!)^S}{(n!)^{DL}},
 &R_n(t)&=\frac{C_nP_n(t)}{\prod_{h=0}^n(t+4h)^S}.
 \label{eq:R}
\end{align}
The zero set of $P_n$ consists of the required grid points in
$[-(r+4)n,rn]$. All these zeros are simple. In particular, one zero
cancels at every denominator point $0,-4,\ldots,-4n$, leaving pole
order at most $s$.

Reflection about the midpoint $-2n$ preserves the zero set.
The numerator has odd degree, so
$P_n(-4n-t)=-P_n(t)$. Since $n$ is odd, the $n+1$ denominator
factors give no sign change. Hence
\[
 R_n(-4n-t)=-R_n(t).
\]
We will multiply $R_n$ by an integer polynomial $W_n$ that is also
odd about $-2n$. The product $T_n=W_nR_n$ will then be reflection
invariant. This symmetry is what removes the odd beta coefficients.
The polynomial $W_n$ will have a fixed homogeneous degree $E<s$
in the variables $n,t$.

\subsection{The criterion to be proved}

To prove decay, we estimate the largest terms in the alternating
sums of $T_n$. The rational number $u$ below bounds where their
exponential size is largest, with distance from the last zero $rn$
measured in units of $n$. Once that location is bounded, the quantity
$U$ bounds the fourth power of the exponential rate after clearing
denominators. Thus the two inequalities locate the maximum and
make the cleared forms tend to zero. The bound $\ep<68/25$ keeps
the criterion rational.

\begin{proposition}\label{prop:criterion}
With the preceding notation, suppose that a positive rational $u$
satisfies
\begin{align}
 f(u)&=\left(\frac{u+L}{u}\right)^{4D}
          \left(\frac{u+r}{u+r+4}\right)^S<1,
          \label{eq:f-criterion}\\
 U&=\left(\frac{68}{25}\right)^{4S}K^4(L+u)^{4DL}4^{4S}
          \left(\frac{r^r}{(r+4)^{r+4}}\right)^S<1.
          \label{eq:U-criterion}
\end{align}
Then $I(s)\ge B$.
\end{proposition}

Section~\ref{sec:gram} constructs the seed constant vector,
Sections~\ref{sec:arithmetic}--
\ref{sec:dilation} prove the arithmetic and analytic estimates, and
Section~\ref{sec:contradiction} completes the contradiction. In
Section~\ref{sec:endpoints} we choose parameters satisfying these
inequalities uniformly in $B$.
